\documentclass[11pt,a4paper]{article}
\usepackage[italian]{babel}
\usepackage[margin=2.5cm]{geometry}
\usepackage{parskip}
\usepackage{amsmath, amssymb, amsthm}
\usepackage{booktabs}
\usepackage{array}
\usepackage{xcolor}
\usepackage{needspace}
\usepackage{enumitem}
\usepackage{listings}
\usepackage{tikz}
\usepackage[most]{tcolorbox}
\usepackage[hidelinks]{hyperref}
\usetikzlibrary{decorations.pathreplacing, shapes.geometric, arrows.meta, positioning}

% Niente vedove/orfane: i paragrafi non lasciano righe isolate a fine/inizio pagina
\widowpenalty=10000
\clubpenalty=10000

% Elenchi più compatti
\setlist{itemsep=2pt, parsep=0pt, topsep=4pt, partopsep=0pt}

% --- Riquadri con bordo nero, senza numero (si spezzano tra due pagine) ---
% Uso: \begin{definizione} ... \end{definizione}
%      \begin{teorema}[Nome] ... \end{teorema}  (il nome è facoltativo)
\tcbset{nota/.style={enhanced jigsaw, breakable, colback=white,
  colframe=black, boxrule=0.5pt, sharp corners}}
\NewTColorBox{definizione}{}{nota,
  before upper={\textbf{Definizione.}\ }}
\NewTColorBox{teorema}{o}{nota,
  before upper={\textbf{Teorema\IfValueT{#1}{ (#1)}.}\ }}
\NewTColorBox{proposizione}{o}{nota,
  before upper={\textbf{Proposizione\IfValueT{#1}{ (#1)}.}\ }}
\NewTColorBox{proprieta}{o}{nota,
  before upper={\textbf{Proprietà\IfValueT{#1}{ (#1)}.}\ }}
\NewTColorBox{assioma}{o}{nota, boxrule=1pt,
  before upper={\textbf{Assioma\IfValueT{#1}{ (#1)}.}\ }}

% --- Ambienti semplici (senza riquadro) ---
% Il titolo sta in cima al blocco, anche se il contenuto inizia con un riquadro o
% con codice; e non resta mai da solo in fondo a una pagina.
% Uso: \begin{esempio}[Titolo facoltativo] ... \end{esempio}
\newcommand{\newrunin}[2]{%
  \NewDocumentEnvironment{#1}{o}{%
    \par\Needspace{4\baselineskip}\addvspace{\medskipamount}\noindent
    \textbf{#2}\IfValueT{##1}{ (##1)}\textbf{.}\ \ignorespaces}%
  {\par\addvspace{\medskipamount}}}
\newrunin{esempio}{Esempio}
\newrunin{esercizio}{Esercizio}
\NewTColorBox{osservazione}{}{nota,
  before upper={\textbf{Osservazione.}\ }}

% V e F nelle tavole di verità
\newcommand{\V}{\textsf{V}}
\newcommand{\F}{\textsf{F}}

% --- Codice C ---
% Uso: \begin{codice} ... \end{codice}      codice C numerato
%      \begin{comando} ... \end{comando}    comandi da terminale
%      \begin{risultato} ... \end{risultato} output di un programma
%      \lstinline|printf("ciao");|           codice dentro al testo
\lstdefinestyle{c}{
  language=C,
  basicstyle=\ttfamily\small,
  keywordstyle=\color{blue}\bfseries,
  commentstyle=\color{gray}\itshape,
  stringstyle=\color{red!70!black},
  numbers=left,
  numberstyle=\tiny\color{gray},
  numbersep=8pt,
  columns=flexible,
  keepspaces=true,
  breaklines=true,
  showstringspaces=false,
  tabsize=4,
  morekeywords={include, define},
  % lettere accentate nei commenti
  literate={à}{{\`a}}1 {è}{{\`e}}1 {é}{{\'e}}1 {ì}{{\`i}}1
           {ò}{{\`o}}1 {ù}{{\`u}}1 {È}{{\`E}}1 {À}{{\`A}}1
}
\lstset{style=c}
\newtcblisting{codice}{listing only, nota, left=6mm,
  listing options={style=c}}
\newtcblisting{comando}{listing only, nota,
  listing options={basicstyle=\ttfamily\small, numbers=none, language={},
    columns=flexible, keepspaces=true, breaklines=true}}
\newtcblisting{risultato}{listing only, enhanced jigsaw, breakable,
  colback=black!5, colframe=black, boxrule=0.5pt, sharp corners,
  listing options={basicstyle=\ttfamily\small, numbers=none, language={},
    columns=flexible, keepspaces=true, breaklines=true}}

% --- Diagrammi di flusso (simboli standard) ---
\tikzset{
  terminale/.style = {ellipse, draw, minimum width=2.5cm, minimum height=0.9cm},
  io/.style        = {trapezium, trapezium left angle=70, trapezium right angle=110,
                      draw, minimum width=2.5cm, minimum height=0.9cm},
  processo/.style  = {rectangle, draw, minimum width=2.5cm, minimum height=0.9cm},
  decisione/.style = {diamond, draw, aspect=2, minimum width=2.5cm},
  freccia/.style   = {thick, -{Stealth}}
}

\title{Conversioni di tipo}
\author{Marco Cerbella}
\date{Lezione 10 -- 13 ottobre 2026}

\begin{document}
\maketitle

\section{Conversioni}

Operandi di tipo diverso si possono combinare in una stessa operazione.

\begin{codice}
double dVar = 2.5;
dVar = dVar * 3;
if (dVar < 10L) { /* ... */ }
\end{codice}

Quando gli operandi hanno tipi diversi, il compilatore cerca di convertirli a un tipo uniforme prima di eseguire l'operazione.

\subsection{Conversioni implicite}

Il compilatore applica conversioni di tipo implicite quando
\begin{itemize}
    \item gli operandi hanno tipi non corrispondenti;
    \item si chiama una funzione con un argomento il cui tipo non corrisponde a quello del parametro;
    \item si inizializzano variabili o si assegnano loro valori.
\end{itemize}
Se la conversione necessaria non è possibile il compilatore segnala un errore; altre volte segnala un warning.

\subsection{Operatore di cast}

Si può convertire un valore da un tipo a un altro anche \emph{esplicitamente}, con l'operatore di cast:
\[
\texttt{(nome\_tipo) espressione}
\]

\begin{codice}
int sum = 10, count = 3;
double mean = (double) sum / count;   // 3.3333
\end{codice}

Il valore di \lstinline|sum| viene prima convertito (esplicitamente) in \lstinline|double|; il compilatore deve poi convertire implicitamente il divisore, il valore di \lstinline|count|, allo stesso tipo prima di eseguire la divisione.

\begin{osservazione}
Conviene usare sempre il cast quando c'è la possibilità di perdere informazione. I cast espliciti evitano i warning del compilatore.
\end{osservazione}

\section{Conversione dei tipi aritmetici}

\subsection{Gerarchia dei tipi}

Quando gli operandi aritmetici hanno tipi diversi, la conversione implicita è governata dal \emph{rango di conversione} dei tipi:
\begin{itemize}
    \item due tipi interi senza segno qualsiasi hanno ranghi diversi: se uno è più largo dell'altro ha rango più alto;
    \item ogni tipo intero con segno ha lo stesso rango del corrispondente tipo senza segno;
    \item i tipi interi standard sono ordinati: \lstinline|_Bool| $<$ \lstinline|char| $<$ \lstinline|short| $<$ \lstinline|int| $<$ \lstinline|long| $<$ \lstinline|long long|;
    \item i tipi in virgola mobile sono ordinati: \lstinline|float| $<$ \lstinline|double| $<$ \lstinline|long double|;
    \item il tipo in virgola mobile di rango più basso, \lstinline|float|, ha rango più alto di qualunque tipo intero;
    \item gli \lstinline|enum| hanno lo stesso rango di \lstinline|int|.
\end{itemize}

\subsection{Promozione intera}

In qualunque espressione si può usare un valore di tipo con rango inferiore a \lstinline|int| al posto di un operando di tipo \lstinline|int| o \lstinline|unsigned int|. In questi casi il compilatore applica la \emph{promozione intera}: ogni operando di rango inferiore a \lstinline|int| viene convertito automaticamente in \lstinline|int|, purché \lstinline|int| sia in grado di rappresentare tutti i valori del tipo originale; altrimenti viene convertito in \lstinline|unsigned int|. Le operazioni nella CPU vengono eseguite su almeno 4 byte.

\begin{esempio}
\begin{codice}
#include <stdio.h>
int main()
{
    char a = 30, b = 40, c = 10;
    char d = (a * b) / c;
    printf("%d ", d);
    return 0;
}
\end{codice}
\begin{risultato}
120
\end{risultato}
A prima vista \lstinline|(a*b)/c| sembra causare un overflow aritmetico: i \lstinline|char| con segno vanno da $-128$ a $127$ (nella maggior parte dei compilatori) e il valore di \lstinline|(a*b)| è $1200$. Ma sui tipi \lstinline|char| avviene la promozione intera, quindi si ottiene il risultato corretto senza overflow.
\end{esempio}

\subsection{Conversioni aritmetiche usuali}

Le \emph{conversioni aritmetiche usuali} sono le conversioni implicite applicate automaticamente agli operandi di tipo aritmetico diverso per la maggior parte degli operatori. Si applicano implicitamente a
\begin{itemize}
    \item gli operatori aritmetici a due operandi: \lstinline|*|, \lstinline|/|, \lstinline|%|, \lstinline|+| e \lstinline|-|;
    \item gli operatori relazionali e di uguaglianza: \lstinline|<|, \lstinline|<=|, \lstinline|>|, \lstinline|>=|, \lstinline|==| e \lstinline|!=|.
\end{itemize}

\begin{esempio}
\begin{codice}
int x = 0;
int i = -1;
unsigned int limit = 200U;
long n = 30L;

if (i < limit)
    x = limit * n;

printf("%d\n", x);
\end{codice}
\begin{risultato}
0
\end{risultato}
\end{esempio}

Le conversioni aritmetiche usuali si applicano così:
\begin{itemize}
    \item Se uno degli operandi ha tipo in virgola mobile, l'operando di rango più basso viene convertito a un tipo con lo stesso rango dell'altro. I tipi reali si convertono solo in tipi reali.
    \item Se entrambi gli operandi sono interi, si esegue prima la promozione intera su entrambi. Se dopo la promozione i tipi sono ancora diversi, si procede così:
    \begin{itemize}
        \item \textbf{Regola 1}: se un operando ha un tipo \emph{senza segno} $T$ con rango almeno pari a quello del tipo dell'altro operando, l'altro operando viene convertito al tipo senza segno $T$.
        \item \textbf{Regola 2}: altrimenti, un operando ha un tipo \emph{con segno} $T$ con rango più alto di quello dell'altro. L'altro operando viene convertito al tipo con segno $T$ solo se $T$ è in grado di rappresentare tutti i valori del suo tipo precedente; se non lo è, entrambi gli operandi vengono convertiti al tipo senza segno corrispondente a $T$.
    \end{itemize}
\end{itemize}

\begin{esempio}[Applicazione delle regole]
Nell'esempio sopra, per valutare il confronto \lstinline|i < limit| il valore di \lstinline|i|, $-1$, viene convertito a \lstinline|unsigned int| (Regola~1: \lstinline|limit| è senza segno e ha rango almeno pari a \lstinline|int|). Il risultato è un numero positivo molto grande, quindi la condizione dell'\lstinline|if| è \emph{falsa} e \lstinline|x| resta $0$.

In \lstinline|limit * n| (Regola~2) \lstinline|limit| viene convertito al tipo di \lstinline|n|, \lstinline|long|, se l'intervallo dei valori di \lstinline|long| contiene tutto l'intervallo di \lstinline|unsigned int|. Altrimenti (ad esempio se \lstinline|int| e \lstinline|long| hanno entrambi $32$ bit) entrambi i fattori vengono convertiti a \lstinline|unsigned long|.

Come si rappresenta $-1$ in un \lstinline|int| (little endian, complemento a due)? Con $32$ bit tutti a $1$: \texttt{11111111111111111111111111111111}. Un \lstinline|unsigned int| ha $32$ bit (in quel compilatore), quindi $-1$ viene convertito implicitamente in \lstinline|unsigned int| e il suo valore diventa $4\,294\,967\,295$.
\end{esempio}

\section{Altre conversioni implicite}

Il compilatore converte automaticamente i valori aritmetici anche nei casi seguenti:
\begin{itemize}
    \item nelle assegnazioni e nelle inizializzazioni, il valore dell'operando di destra viene sempre convertito al tipo dell'operando di sinistra;
    \item nelle chiamate di funzione, gli argomenti vengono convertiti ai tipi dei parametri corrispondenti;
    \item nelle istruzioni \lstinline|return|, il valore dell'espressione viene convertito al tipo restituito dalla funzione.
\end{itemize}

In un'\emph{assegnazione composta}, come \lstinline|x += 2.5| (cioè \lstinline|x = x + 2.5|), i valori di entrambi gli operandi sono prima soggetti alle conversioni aritmetiche usuali; poi il risultato dell'operazione viene convertito, come in una assegnazione semplice, al tipo dell'operando di sinistra. Se \lstinline|x| è \lstinline|int|, viene convertito in \lstinline|double| e poi il risultato \lstinline|x + 2.5| (di tipo \lstinline|double|) viene riconvertito in \lstinline|int|.

\begin{esempio}
\begin{codice}
#include <math.h>   // dichiara double sqrt(double)
int i = 7;
float x = 0.5;      // la costante viene convertita da double a float

i = x;              // il valore di x viene convertito da float a int

x += 2.5;           // prima dell'addizione x viene convertito in double;
                    // la somma viene poi convertita in float per l'assegnazione

x = sqrt(i);        // l'argomento viene convertito da int a double;
                    // il valore restituito da double a float

long my_func() {    // la costante 0 viene convertita a long,
    /* ... */       // il tipo restituito dalla funzione
    return 0;
}
\end{codice}
\end{esempio}

\section{Conversioni verso tipi interi}

\subsection{Verso tipi interi senza segno}

I valori interi sono sempre preservati se rientrano nell'intervallo del nuovo tipo senza segno, cioè tra $0$ e \lstinline|Utype_MAX|. Per i valori fuori intervallo, il valore dopo la conversione si ottiene sommando o sottraendo $\mathtt{Utype\_MAX} + 1$ tante volte quante servono perché il risultato rientri nell'intervallo del nuovo tipo.

\begin{codice}
unsigned short n = 1000;   // 1000 rientra nell'intervallo di unsigned short: ok
n = -1;                    // il valore -1 deve essere convertito
\end{codice}

Si ha $-1 + (\mathtt{USHRT\_MAX} + 1) = \mathtt{USHRT\_MAX}$, quindi l'ultima istruzione equivale a \lstinline|n = USHRT_MAX;|.

\subsection{Da virgola mobile a interi}

Per convertire un numero reale in virgola mobile in un tipo intero (con o senza segno) il compilatore scarta la parte frazionaria. Se la parte intera rimanente è fuori dall'intervallo del nuovo tipo, il risultato della conversione è \emph{indefinito}.

\begin{codice}
double x = 2.9;
unsigned long n = x;   // la parte frazionaria di x viene persa: n = 2
\end{codice}

\subsection{Verso tipi interi con segno}

Il problema del superamento dell'intervallo dei valori del tipo di destinazione può presentarsi anche quando un valore intero, con o senza segno, viene convertito in un tipo intero \emph{con segno}, ad esempio da \lstinline|long| o \lstinline|unsigned int| a \lstinline|int|. A differenza delle conversioni verso tipi senza segno, il risultato di un overflow in conversione verso un tipo con segno è lasciato all'implementazione. La maggior parte dei compilatori scarta i bit più significativi della rappresentazione binaria del valore originale e interpreta i bit più bassi secondo il nuovo tipo.

\begin{esempio}
\begin{codice}
#include <stdio.h>
#include <limits.h>
int main() {
    long long int a = (LLONG_MAX - UINT_MAX) + 1;
    int b = a;

    printf("b: %d\n", b);
    printf("a: %lld\n", a);
}
\end{codice}
\begin{risultato}
b: 1
a: 9223372032559808513
\end{risultato}
In esadecimale \lstinline|LLONG_MAX| è \texttt{0x7FFFFFFFFFFFFFFF} e \lstinline|UINT_MAX| è \texttt{0xFFFFFFFF}, quindi
\[
a = \texttt{0x7FFFFFFF00000001}.
\]
Convertendo in \lstinline|int| (32 bit) si tengono solo i $32$ bit più bassi, \texttt{00000001}: perciò $b = 1$.
\end{esempio}

\section{Conversioni verso tipi in virgola mobile}

Non tutti i valori interi sono rappresentabili esattamente nei tipi in virgola mobile. Ad esempio, benché l'intervallo di valori del tipo \lstinline|float| includa quello di \lstinline|long| e \lstinline|long long|, \lstinline|float| ha una precisione di sole sei cifre decimali. Quindi alcuni valori \lstinline|long| non possono essere memorizzati esattamente in un \lstinline|float|: il risultato della conversione è il valore rappresentabile immediatamente inferiore o superiore.

\begin{esempio}
\begin{codice}
float r_var = 16777217;
double l_var = 16777217;
printf("The rounding error (l_var - r_var) is %.2f\n", l_var - r_var);
printf("r_var is %.2f\n", r_var);
\end{codice}
\begin{risultato}
The rounding error (l_var - r_var) is 1.00
r_var is 16777216.00
\end{risultato}
\end{esempio}

Ogni valore di un tipo in virgola mobile si può rappresentare esattamente in un altro tipo in virgola mobile di precisione maggiore. Quindi, quando un \lstinline|double| viene convertito in \lstinline|long double|, o un \lstinline|float| in \lstinline|double| o \lstinline|long double|, il valore è esattamente preservato.

Nelle conversioni da un tipo più preciso a uno meno preciso, invece, il valore da convertire può essere fuori dall'intervallo del nuovo tipo.
\begin{itemize}
    \item Se il valore supera l'intervallo del tipo di destinazione, il risultato della conversione è \emph{indefinito}.
    \item Se il valore è nell'intervallo ma non è rappresentabile esattamente con la precisione del tipo di destinazione, il risultato è il valore rappresentabile immediatamente inferiore o superiore.
\end{itemize}

\section*{Riferimenti}

Niente sul libro. Risorse indicate dal docente:
\begin{itemize}
    \item \url{https://www.w3schools.com/c/c_type_conversion.php}
    \item \url{https://www.guru99.com/c-type-casting.html}
    \item soprattutto, il capitolo 4 di \emph{C in a Nutshell} (O'Reilly):\\
          \url{https://www.oreilly.com/library/view/c-in-a/0596006977/ch04.html}
\end{itemize}

\end{document}
